\section{Arithmetic coefficient maps and their singular values}
\label{am-section}

The coefficient lattices of Section~\ref{coef:section} give explicit
arithmetic maps whose image ideals and archimedean norms can be computed.
This section proves the estimates needed for these maps, including the
normalizing constant in their metrics. It also gives sharper bounds in
terms of elementary divisors and singular values. The comparison with the
actual class components is stated at the end. These are coefficient maps;
the different valuation inequalities for the iterated period pairing
require additional arguments.\cite{alpo-uniform-source}

\subsection{The actual map and its integral image}

Let $K$ be a number field, $\mathcal O$ its ring of integers, and $n\ge1$.
Fix a finite set $\mathcal T$. For each $\tau\in\mathcal T$ let
$\lambda_\tau=(\lambda_{\tau,1},\ldots,\lambda_{\tau,n})$ be a
nonincreasing tuple of integers. Put
\[
 a_\tau=\lambda_{\tau,n},\qquad
 \nu_{\tau,i}=\lambda_{\tau,i}-a_\tau,\qquad
 d_\tau=\sum_i\nu_{\tau,i}.
\]
Thus $\nu_\tau$ is a partition and $\nu_{\tau,n}=0$.
The arguments below also apply to a partition with a nonzero last part
and an independently specified integer $a_\tau$.

Use a separate $n$ by $n$ variable matrix $X_\tau$ for each $\tau$ and set
\[
 V=\bigotimes_{\tau\in\mathcal T}V_{\nu_\tau,K},\qquad
 \Delta_0=\bigotimes_\tau\Delta_{\nu_\tau}(\mathcal O),\qquad
 \nabla_0=\bigotimes_\tau\nabla_{\nu_\tau}(\mathcal O).
\]
The tensor products are realized as products of polynomials in the
disjoint variable blocks. Define
\[
 \rho_0(g)F((X_\tau)_\tau)=F((g_\tau^tX_\tau)_\tau),\qquad
 \rho_\lambda(g)=\left(\prod_\tau\det(g_\tau)^{a_\tau}\right)\rho_0(g).
 \tag{AM1}\label{am-action}
\]
Both are representations of $\prod_\tau\operatorname{GL}_n(K)$:
substitution reverses multiplication, and the transpose reverses it
again. The determinant powers in the second expression may be negative.

Let $\mathcal A$ be the finite set of arrays
$\alpha=(\alpha_{\tau,i,s})$ of nonnegative integers satisfying
\[
             \sum_{i=1}^n\alpha_{\tau,i,s}=\nu_{\tau,s}
                       \quad\text{for every $\tau,s$}.
\]
Every polynomial in $V$ is supported on $\mathcal A$.
For $g_\tau\in\operatorname{GL}_n(K)$ the coefficient map is
\begin{equation}\label{am-map}
 C_g:V\longrightarrow K^{\mathcal A},\qquad
 C_g(F)_\alpha=[X^\alpha]\rho_0(g)F.
\end{equation}
In particular the determinant character is \emph{not} inserted into
the definition of $C_g$. It will occur in its integral image ideal.
The map is injective: all the coefficients of $\rho_0(g)F$ occur in
\eqref{am-map}, and $\rho_0(g^{-1})$ is its inverse substitution.

The translates of $\prod_\tau H_{\nu_\tau}$ under $\rho_\lambda$ span
$V$. Here is a direct justification of the spanning statement. The
integral duality theorem proves that the divided-row-operation words
span $V$ over $K$. Each such word is a coefficient of a polynomial in
finitely many shear parameters. Polynomial interpolation recovers that
coefficient from finitely many values at nonnegative integers. Every
off-diagonal shear is invertible, and a diagonal substitution
$1+t$ is also invertible at these parameter values. Thus the relevant
values are invertible-matrix translates of the extremal polynomial.
Their determinant characters are nonzero scalars, which do not change
their $K$-span. Tensoring the argument proves the assertion for all blocks.

For full $\mathcal O$-lattices $M_\tau\subset K^n$, write
$\mathfrak d(M_\tau)\subset K$ for their determinant ideals, using
$e_1\wedge\cdots\wedge e_n$ to identify $\bigwedge^nK^n$ with $K$.
At a nonzero prime $\mathfrak p$, choose local bases
\[
       M_{\tau,\mathfrak p}=B_{\tau,\mathfrak p}\mathcal O_{\mathfrak p}^n.
\]
Define the transported coefficient lattice $\Delta_\lambda(M)$ by
\begin{equation}\label{am-local-lattice}
 \Delta_\lambda(M)_{\mathfrak p}
       =\rho_\lambda(B_{\mathfrak p})\Delta_{0,\mathfrak p}.
\end{equation}
The lattice $\nabla_\lambda(M)$ is defined in the same way with
$\nabla_0$. These definitions agree with the intersection over all
local bases used in Section~\ref{coef:section}.

To check independence of the bases, column invariance immediately
shows that $\rho_0(U)$ preserves $\nabla_0$ when
$U_\tau\in\operatorname{GL}_n(\mathcal O_{\mathfrak p})$.
It preserves $\Delta_0$ as well: the integral duality already proved in
Theorem~\ref{coef:integral-duality} and the transpose adjoint identity
identify $\Delta_0$ with the dual of $\nabla_0$ under the normalized
bilinear Fock pairing. Apply the adjoint identity also to $U^{-1}$ to
obtain equality. The determinant characters are units on this group,
even for negative $a_\tau$. Only finitely many local lattices in
\eqref{am-local-lattice} differ from the standard one, so their
intersection is a full lattice, by the elementary local patching result
of Section~\ref{div:section}.

\begin{theorem}[Integral coefficient map]\label{am-integral}
Suppose $g_\tau M_\tau\subset\mathcal O^n$ for every $\tau$, and put
\begin{equation}\label{am-ideal}
                  \mathfrak I=\prod_\tau\mathfrak d(M_\tau)^{a_\tau}.
\end{equation}
Then $\mathfrak I$ is a nonzero fractional ideal and
\[
 C_g(\Delta_\lambda(M))\subset\mathfrak I^{\mathcal A},
 \qquad
 C_g(\nabla_\lambda(M))\subset\mathfrak I^{\mathcal A}.
 \tag{AM2}\label{am-image}
\]
Equivalently, every coordinate of the image has valuation at least
$\sum_\tau a_\tau\operatorname{ord}_{\mathfrak p}\det B_{\tau,\mathfrak p}$
at each $\mathfrak p$. The assertion is valid at primes above $2$ and
does not require any $M_\tau$ to be free over $\mathcal O$.
\end{theorem}
\begin{proof}
Localize at $\mathfrak p$ and write
$F=\rho_\lambda(B)F_0$, where $F_0$ belongs to the appropriate standard
lattice. Its coefficients are integral over $\mathcal O_{\mathfrak p}$.
Since $g_\tau B_\tau$ has integral entries,
\[
 \rho_0(g)F
   =\left(\prod_\tau\det B_\tau^{a_\tau}\right)\rho_0(gB)F_0
\]
has every coefficient in
$\prod_\tau\det B_\tau^{a_\tau}\mathcal O_{\mathfrak p}$.
This is exactly $\mathfrak I_{\mathfrak p}$, because
$\mathfrak d(M_\tau)_{\mathfrak p}=\det B_\tau\mathcal O_{\mathfrak p}$.
Intersect these inclusions over $\mathfrak p$. The same description also
proves that the intersection defined by the conditions
\[
 \left(\prod_\tau\det B_\tau^{-a_\tau}\right)c
                              \in\mathcal O_{\mathfrak p}
\]
for all local bases is precisely the fractional ideal \eqref{am-ideal}.
\end{proof}

\subsection{Row weights and a sharper local divisibility statement}

For a polynomial in one variable block, its row degree is
$\beta=(\beta_1,\ldots,\beta_n)$ if every monomial has total degree
$\beta_i$ in row $i$. Projection to a row degree preserves the defining
column identities of $V_\nu$.

\begin{lemma}[Row weights]\label{am-weights}
If the row-degree-$\beta$ subspace of $V_\nu$ is nonzero, then
$\sum_i\beta_i=\sum_i\nu_i$, and for every subset
$S\subset\{1,\ldots,n\}$ with $|S|=k$,
\begin{equation}\label{am-majorization}
       \sum_{i\in S}\beta_i\le\sum_{i=1}^k\nu_i.
\end{equation}
In particular, for every $k$,
\[
 \sum_{i=n-k+1}^n\nu_i\le\sum_{i=1}^k\beta_i
                              \le\sum_{i=1}^k\nu_i.
\]
The weights $\nu$ and $(\nu_n,\ldots,\nu_1)$ both occur.
\end{lemma}
\begin{proof}
For $S=\{1,\ldots,k\}$, restrict a nonzero polynomial of row degree
$\beta$ to lower triangular matrices. The restriction is nonzero by
Lemma~\ref{coef:restriction}. In a monomial of this restriction, the
first $k$ rows use only the first $k$ columns. Their total degree is
therefore at most $\sum_{i\le k}\nu_i$. Row permutations preserve
column invariance and column degrees, which proves the assertion for
every $S$. The lower bound follows by applying the upper bound to the
complement of the first $k$ rows. The extremal polynomial $H_\nu$ has
row degree $\nu$; reverse its rows for the other endpoint.
\end{proof}

The lemma requires the column-invariance identities. Column homogeneity
alone would be insufficient. For example, in three rows the polynomial
$x_{11}x_{12}$ has column degrees $(1,1,0)$ and row degree $(2,0,0)$;
it violates \eqref{am-majorization} and is not column invariant.

\begin{theorem}[Exact local coefficient ideal]\label{am-smith}
Let $R$ be a discrete valuation ring of characteristic zero with fraction
field $L$, uniformizer $\pi$, and normalized valuation $v(\pi)=1$.
Let $M_\tau=B_\tau R^n$ and $g_\tau\in\operatorname{GL}_n(L)$, without
requiring $g_\tau M_\tau\subset R^n$. Let
$c_{\tau,1}\le\cdots\le c_{\tau,n}$ be the elementary-divisor
exponents of $g_\tau B_\tau$. The fractional ideal generated by all
coordinates $C_g(F)_\alpha$, as $F$ ranges over
$\Delta_\lambda(M)$, is exactly
\begin{equation}\label{am-sharp-local}
 \pi^eR,\qquad
 e=\sum_\tau\left(a_\tau v(\det B_\tau)
                         +\sum_i\nu_{\tau,i}c_{\tau,i}\right).
\end{equation}
The same statement holds with $\nabla_\lambda(M)$.
\end{theorem}
\begin{proof}
Smith reduction over $R$ gives
$g_\tau B_\tau=U_\tau D_\tau V_\tau$ with $U_\tau,V_\tau$ unimodular
and $D_\tau=\operatorname{diag}(\pi^{c_{\tau,1}},\ldots,
\pi^{c_{\tau,n}})$. This also holds for matrices with denominators:
first multiply the matrix by a power of $\pi$ and then subtract that
power from every elementary-divisor exponent.
For completeness, integral Smith reduction follows by moving an entry
of least valuation to the upper left. It divides every entry. Integral
row and column subtractions clear its row and column; induction treats
the remaining block, whose entries are still divisible by the chosen
pivot. Thus the successive pivot valuations are nondecreasing.

Integral unimodular substitutions preserve both standard lattices and
preserve the ideal generated by the coefficients of any polynomial:
one inclusion follows from integrality of the substitution, and the
reverse inclusion follows by applying its integral inverse.
It remains to treat the diagonal matrices $D_\tau$.
On a row-degree-$\beta$ term, $\rho_0(D_\tau)$ multiplies by
$\pi^{\sum_i c_{\tau,i}\beta_i}$. Summation by parts and
Lemma~\ref{am-weights} give
\[
 \sum_i c_{\tau,i}\beta_i
 =c_{\tau,n}d_\tau
   -\sum_{k<n}(c_{\tau,k+1}-c_{\tau,k})\sum_{i\le k}\beta_i
 \ge\sum_i c_{\tau,i}\nu_{\tau,i}.
\]
The determinant character of $B$ supplies the first term of $e$.
For equality take the product of the $H_{\nu_\tau}$ after undoing
the substitutions $V_\tau$. Each $H_{\nu_\tau}$ belongs to the Weyl
lattice, has row degree $\nu_\tau$, and has a coefficient equal to one:
its restriction to lower triangular matrices is the diagonal monomial
with coefficient one. Hence the lower bound is attained as the ideal
of coefficients, even after the substitutions $U_\tau$.
\end{proof}

In Theorem~\ref{am-integral} all $c_{\tau,i}$ are nonnegative, so
\eqref{am-sharp-local} implies the stated determinant-ideal bound.
The sharper formula also identifies exactly how much divisibility is
lost when the chosen linear functionals have denominators.

\subsection{The Fock metric and exact singular-value bounds}

Over $\mathbb C$ put
\[
 C_\tau=\sum_\alpha\alpha!\,|[X^\alpha]H_{\nu_\tau}|^2,
 \qquad C=\prod_\tau C_\tau>0,
 \qquad
 \|F\|_V^2=C^{-1}\sum_\alpha\alpha!\,|F_\alpha|^2.
\]
Give $\mathbb C^{\mathcal A}$ the diagonal metric
\begin{equation}\label{am-coordinate-metric}
 \|z\|_{\mathcal A}^2
       =\sum_{\alpha\in\mathcal A}w_\alpha|z_\alpha|^2,
       \qquad w_\alpha=\alpha!/C.
\end{equation}
Then $C_1$ is an isometric embedding. This explains the factorial weights
in the target metric; replacing them by the ordinary coordinate metric
would change the operator estimates.

\begin{theorem}[Largest and smallest singular values]\label{am-singular}
For each $\tau$ let
$s_{\tau,1}\ge\cdots\ge s_{\tau,n}>0$ be the singular values of
$g_\tau\in\operatorname{GL}_n(\mathbb C)$. Then
\begin{align}
 \|C_g\|_{\mathrm{op}}
       &=\prod_{\tau,i}s_{\tau,i}^{\nu_{\tau,i}},
       \label{am-op-exact}\\
 \inf_{F\ne0}\frac{\|C_gF\|_{\mathcal A}}{\|F\|_V}
       &=\prod_{\tau,i}s_{\tau,i}^{\nu_{\tau,n+1-i}}.
       \label{am-min-exact}
\end{align}
For $\rho_\lambda(g)$ multiply both expressions by
$\prod_\tau|\det g_\tau|^{a_\tau}$; equivalently replace $\nu$ by
$\lambda$ in the corresponding ordered products.
In particular
\begin{equation}\label{am-frobenius}
 \|C_gF\|_{\mathcal A}^2
 \le\left(\prod_\tau\|g_\tau\|_{\mathrm{op}}^{2d_\tau}\right)\|F\|_V^2
 \le\left(\prod_\tau\|g_\tau\|_{\mathrm{HS}}^{2d_\tau}\right)\|F\|_V^2.
\end{equation}
\end{theorem}
\begin{proof}
The Hermitian version of Lemma~\ref{coef:adjoint} says that the adjoint
of $\rho_0(g)$ for the Fock metric is $\rho_0(g^*)$.
Thus unitary substitutions act unitarily on $V$. A singular-value
decomposition reduces the problem to positive diagonal $g_\tau$.
Distinct row weights are orthogonal, and on the weight $\beta$ the
diagonal action has eigenvalue $\prod_i s_{\tau,i}^{\beta_i}$.
Writing $u_i=\log s_{\tau,i}$, summation by parts gives
\[
 \sum_i u_i\beta_i=u_nd_\tau
                   +\sum_{k<n}(u_k-u_{k+1})\sum_{i\le k}\beta_i.
\]
The two partial-sum bounds of Lemma~\ref{am-weights} prove the upper
and lower estimates. Its endpoint weights attain both estimates.
Their tensor products attain the endpoints simultaneously for all
$\tau$. Finally $C_1$ is isometric, so the singular values of $C_g$
are those of $\rho_0(g)$ on $V$. The scalar determinant twist gives
the last assertion. The two inequalities in \eqref{am-frobenius}
follow from $s_{\tau,i}\le s_{\tau,1}\le\|g_\tau\|_{\mathrm{HS}}$.
\end{proof}

No bound involving only $\det g$ could replace
\eqref{am-op-exact}: for $\nu=(1,0)$ and
$g=\operatorname{diag}(T,T^{-1})$, $T>1$, the determinant is one
and the operator norm is $T$. Also, the column-homogeneous polynomial
$x_{11}x_{12}$ above is multiplied by $T^2$ under
$\operatorname{diag}(T,1,1)$, whereas the claimed bound for partition
$(1,1,0)$ would be $T$. This confirms the necessity of column
invariance in the exact formula.

\begin{proposition}[Determinant of the substitution]\label{am-determinant}
Let $r=\dim_KV$. The numbers $\kappa_\tau=rd_\tau/n$ are integers, and
\[
 \det(\rho_0(g)|V)=\prod_\tau\det(g_\tau)^{\kappa_\tau},\qquad
 \det(\rho_\lambda(g)|V)
                 =\prod_\tau\det(g_\tau)^{\kappa_\tau+ra_\tau}.
 \tag{AM3}\label{am-det-formula}
\]
Over $\mathbb C$, in orthonormal source coordinates,
$\det(C_g^*C_g)=\prod_\tau|\det g_\tau|^{2\kappa_\tau}$.
\end{proposition}
\begin{proof}
Decompose $V$ into row-weight spaces in block $\tau$.
Row permutations identify the relevant weight spaces, so the sum of
the $i$th row degrees, counted with their dimensions, is independent
of $i$. Its sum over $i$ is $rd_\tau$; hence each sum is the integer
$\kappa_\tau$. This proves \eqref{am-det-formula} on diagonal matrices.
An elementary shear has unipotent action on the polynomial space of
degree $d_\tau$, since it changes one row degree in only one direction.
Its restriction to the invariant space $V$ has determinant one.
Gaussian elimination generates every invertible matrix by diagonal
matrices and elementary shears (a row interchange is a product of
three shears and a diagonal matrix). Thus the formula holds for every
$g$. The determinant twist acts by a scalar on an $r$-dimensional
space. The last formula follows because $C_1$ is isometric.
\end{proof}

\subsection{An elementary bound for the normalizing constant}

The preceding formulas retain the exact constant $C_\nu$.
Its required upper bound has a short proof by Gaussian integration;
no dimension formula for polynomial representations is needed.

\begin{theorem}[Norm of the extremal polynomial]\label{am-gaussian}
For every partition $\nu$ of length $n\ge1$,
\begin{equation}\label{am-normalization-bound}
 \prod_i\nu_i!\ \le\ C_\nu\ \le\
 \prod_{i=1}^n\frac{(\nu_i+n-i)!}{(n-i)!}
 \ \le\left(\prod_i\nu_i!\right)(\nu_1+n)^{n(n-1)/2}.
\end{equation}
\end{theorem}
\begin{proof}
The diagonal monomial in $H_\nu$ has coefficient one, so its contribution
to the Fock norm is $\prod_i\nu_i!$, proving the lower bound.

Let $Z$ be an $n$ by $n$ matrix of independent complex Gaussian entries,
each with density $\pi^{-1}e^{-|z|^2}$ on $\mathbb C$.
Polar integration gives
\[
 \mathbb E(z^a\overline{z}^{b})=
 \begin{cases}a!&a=b,\\0&a\ne b.\end{cases}
\]
Indeed the angular integral vanishes unless $a=b$; the remaining radial
integral is $\int_0^\infty t^ae^{-t}\,dt=a!$, by integration by parts.
Independence consequently proves
$C_\nu=\mathbb E|H_\nu(Z)|^2$.

Let $Z_j$ consist of the first $j$ columns of $Z$.
Cauchy--Binet gives
\[
           |\Delta_j(Z)|^2\le\det(Z_j^*Z_j).
\]
Apply Gram--Schmidt to these columns, writing $r_i>0$ for the length
of the component of column $i$ orthogonal to the preceding columns.
The columns are independent with probability one: the zero set of a
nonzero polynomial has Lebesgue measure zero, by induction on the
number of variables and Fubini's theorem. Thus
\[
 |H_\nu(Z)|^2
 \le\prod_j\det(Z_j^*Z_j)^{\nu_j-\nu_{j+1}}
 =\prod_i r_i^{2\nu_i},\qquad \nu_{n+1}=0.
\]
Conditionally on the preceding columns, the orthogonal component of
column $i$ is a standard complex Gaussian vector in dimension $n-i+1$.
This follows directly from unitary invariance of its density.
For such a vector in dimension $q$, the multinomial theorem and the
one-variable moments give
\[
 \mathbb E\left(\sum_{j=1}^q|z_j|^2\right)^m
 =\sum_{b_1+\cdots+b_q=m}\frac{m!}{b_1!\cdots b_q!}
                                      \prod_j b_j!
 =\frac{(m+q-1)!}{(q-1)!}.
\]
There are $\binom{m+q-1}{q-1}$ nonnegative $q$-tuples with sum $m$,
as follows by inserting $q-1$ dividers into a string of $m$ symbols.
Successive conditional integrations, starting with the last column,
therefore prove the middle inequality of \eqref{am-normalization-bound}.
Finally
\[
 \frac{(\nu_i+n-i)!}{(n-i)!\nu_i!}
          =\prod_{j=1}^{n-i}\frac{\nu_i+j}{j}
          \le(\nu_1+n)^{n-i}.
\]
Multiply over $i$.
\end{proof}

For $m\ge0$ write $m=nq+r$ with $0\le r<n$, and put
\[
                    b_n(m)=(q!)^{n-r}((q+1)!)^r.
\]
The least factorial weight in \eqref{am-coordinate-metric} is exactly
\begin{equation}\label{am-min-weight}
 w_{\min}=\min_{\alpha\in\mathcal A}w_\alpha
                    =\frac{\prod_{\tau,s}b_n(\nu_{\tau,s})}{C}.
\end{equation}
To prove this, minimize the product of the factorials separately in
each column. If two entries are $u\ge v+2$, replacing them by $u-1,v+1$
decreases the product by the factor $(v+1)/u<1$. Repetition leaves
entries differing by at most one, which are precisely the balanced
entries in the definition of $b_n$.

A less precise bound, but one expressed without factorial sums, is
\begin{equation}\label{am-weight-lower}
 w_{\min}\ge
 \left(n^{\sum_\tau d_\tau}
       \prod_\tau\prod_{i=1}^n
        \frac{(\nu_{\tau,i}+n-i)!}{(n-i)!\nu_{\tau,i}!}\right)^{-1}
 \ge
 \left(n^{\sum_\tau d_\tau}
       \prod_\tau(\nu_{\tau,1}+n)^{n(n-1)/2}\right)^{-1}.
\end{equation}
Indeed the multinomial theorem at $(1,\ldots,1)$ implies
$m!\le n^m\prod_i b_i!$ whenever $\sum_i b_i=m$.
Multiply this over the columns and apply Theorem~\ref{am-gaussian}.

\subsection{The resulting slope bound, with ordinary degree conventions}

Write $d_K=[K:\mathbb Q]$, and let $n_v=1$ or $2$ be the local degree
at a real or complex archimedean place. All metrics below are compatible
with complex conjugation. For a Hermitian lattice $L$ of rank $r>0$,
choose $0\ne\omega\in\bigwedge^r_K(KL)$ and write
$\bigwedge^r_{\mathcal O}L=\mathfrak a\omega$.
Our normalized degree and slope are
\begin{equation}\label{am-degree}
 \widehat{\deg}L=-\frac{\log N\mathfrak a}{d_K}
      -\frac1{d_K}\sum_{v\mid\infty}n_v\log\|\omega\|_v,
 \qquad \mu(L)=\widehat{\deg}L/r.
\end{equation}
Here $N\mathfrak a$ is the absolute norm of a fractional ideal, a
positive rational number. The product formula makes the definition
independent of $\omega$. Let $\mu_{\max}(L)$ be the supremum of the
slopes of its nonzero sublattices, with their induced metrics.

\begin{lemma}[A diagonal target]\label{am-diagonal-slope}
Suppose $\mathfrak I\subset K$ is a nonzero fractional ideal and
$K^A$ has diagonal metrics with constant positive diagonal weights
$w_a\ge w_0>0$ at every infinite place. Then every nonzero sublattice
$L\subset\mathfrak I^A$ satisfies
\[
                 \mu(L)\le-\frac{\log N\mathfrak I}{d_K}
                                      -\frac12\log w_0.
\]
\end{lemma}
\begin{proof}
Let $r$ be its rank, write $\bigwedge^rL=\mathfrak a\omega$, and
choose a nonzero coordinate $b$ of the exterior vector $\omega$.
Since all coordinates of $\bigwedge^rL$ lie in $\mathfrak I^r$,
$b\mathfrak a\subset\mathfrak I^r$ and hence
$\log N\mathfrak a+\log|N_{K/\mathbb Q}b|
\ge r\log N\mathfrak I$. In the exterior metric,
$\|\omega\|_v\ge w_0^{r/2}|v(b)|$, because its squared norm is the
sum of its squared coordinates times products of $r$ diagonal weights.
Insert these two inequalities into \eqref{am-degree} and use
$\sum_vn_v\log|v(b)|=\log|N_{K/\mathbb Q}b|$.
\end{proof}

For an injective $K$-linear map $f$, put
$h(f)=d_K^{-1}\sum_v n_v\log\|f\|_{v,\mathrm{op}}$.
For every sublattice $L$, the exterior operator estimate gives
$\widehat{\deg}L\le\widehat{\deg}f(L)+r h(f)$: use the same
fractional ideal in the determinant generators $\omega$ and $f\omega$,
and $\|f\omega\|_v\le\|f\|_{v,\mathrm{op}}^r\|\omega\|_v$.
This proves the following bound directly for the map constructed here.

\begin{corollary}[Slope of the coefficient lattice]\label{am-slope}
Give $\Delta_\lambda(M)$ the normalized Fock metric at every infinite
place and suppose $g_\tau M_\tau\subset\mathcal O^n$.
With $\mathfrak I$ as in \eqref{am-ideal} and the singular values of
$v(g_\tau)$ ordered decreasingly,
\begin{equation}\label{am-slope-bound}
 \mu_{\max}(\Delta_\lambda(M))
 \le-\frac{\log N\mathfrak I}{d_K}-\frac12\log w_{\min}
       +\frac1{d_K}\sum_v n_v\sum_{\tau,i}
                                  \nu_{\tau,i}\log s_{\tau,i}(v).
\end{equation}
The same bound holds for $\nabla_\lambda(M)$.
One may replace the last summand by
\[
 \frac1{2d_K}\sum_vn_v\sum_\tau d_\tau
                          \log\|v(g_\tau)\|_{\mathrm{HS}}^2
\]
and bound $w_{\min}$ by \eqref{am-weight-lower}.
\end{corollary}
\begin{proof}
Apply Theorem~\ref{am-integral}, Lemma~\ref{am-diagonal-slope}, and the
preceding exterior operator inequality to $C_g$. Its injectivity and
its exact operator norm were proved above. The weaker version is
\eqref{am-frobenius}.
\end{proof}

The use of ordinary slopes in \eqref{am-degree} matters. A convention
using logarithms of squared norms doubles both the slope and the height.
Under that convention the right side of \eqref{am-slope-bound} must
also be doubled; in particular the ideal term becomes
$-2\log N\mathfrak I/d_K$.

\subsection{The maps on actual class components}

We describe precisely the additional data needed to apply the result
to class components, and prove the resulting comparison.
In the application, class values are obtained by evaluating equivariant
functions on representatives of decorated-lattice isometry classes.
At the level in question the stabilizers are trivial. Evaluation then
identifies the space of such functions with the full product of their
coefficient spaces: assign arbitrary values at the representatives and
extend by the coefficient action. The extension is well defined because
two isometries with the same image differ by a stabilizer. Conversely,
equivariance determines the function from these values. The class-sum
metric has reciprocal stabilizer weights, which are one at this level.
For nontrivial stabilizers one must instead use their invariant
coefficient spaces; the full-product assertion would be false in general.

Let $\mathcal C$ be finite. For every $c\in\mathcal C$ and
$\tau\in\mathcal T$, let $M_{c,\tau}\subset K^n$ be a full lattice.
Consider the lattice in $\bigoplus_cV$ whose localization at
$\mathfrak p$ consists of the tuples $(F_c)_c$ satisfying
\begin{equation}\label{am-class-local}
 \rho_\lambda(B_{c,\mathfrak p}^{-1})F_c
                               \in\Delta_{0,\mathfrak p}\quad(c\in\mathcal C),
\end{equation}
where $B_{c,\tau,\mathfrak p}$ are local bases of $M_{c,\tau}$.
This lattice is exactly
$\mathcal L=\bigoplus_c\Delta_\lambda(M_c)$.
Consequently the projection onto any class is integral, and the
projection is injective on its class summand. The same assertions
remain true if the right side of \eqref{am-class-local} is replaced by
the integral span of the same divided-row words: this is the defining
Weyl lattice, as proved in Section~\ref{coef:section}.

At every infinite place choose invertible matrices $X_{c,\tau,v}$ and
give $K^n$ the metric
\[
       h_{c,\tau,v}(z,z')=(X_{c,\tau,v}^{-1}z)^*
                                      (X_{c,\tau,v}^{-1}z').
\]
Give $\mathcal L$ the direct-sum metric
\begin{equation}\label{am-class-metric}
       \|F\|_v^2=\sum_c
             \|\rho_\lambda(X_{c,v}^{-1})F_c\|_V^2.
\end{equation}
Fix a class $c$. Choose, for every $\tau$, a basis
$\ell_{\tau,1},\ldots,\ell_{\tau,n}$ of $(K^n)^*$ which is integral
on $M_{c,\tau}$, and let $g_\tau$ be its row matrix.
Define the actual map
\[
                     \Phi_c(F)=C_g(F_c).
\]
It is injective on the $c$th summand and takes that summand into
$\mathfrak I_c^{\mathcal A}$, where
$\mathfrak I_c=\prod_\tau\mathfrak d(M_{c,\tau})^{a_\tau}$.

At $v$, write $y_{\tau,j,v}$ for the Riesz vector of
$v(\ell_{\tau,j})$ for $h_{c,\tau,v}$, using the convention
$\ell(z)=h(y,z)$, conjugate-linear in the first argument.
Let $R_{\tau,v}$ be the matrix whose $j$th row is
$(X_{c,\tau,v}^{-1}y_{\tau,j,v})^*$. Then
\begin{equation}\label{am-riesz}
       v(g_\tau)X_{c,\tau,v}=R_{\tau,v},\qquad
       \|R_{\tau,v}\|_{\mathrm{HS}}^2
                         =\sum_j\|y_{\tau,j,v}\|_{h_{c,\tau,v}}^2.
\end{equation}
Both identities follow by writing
$h(y,z)=y^*(X^{-1})^*X^{-1}z$.
They retain the conjugates; replacing $y^*$ by $y^t$ would generally
give a different functional.

\begin{proposition}[Class coefficient bound]\label{am-class-bound}
Let $s_{\tau,1}(R_{\tau,v})\ge\cdots\ge s_{\tau,n}(R_{\tau,v})>0$
be the singular values in \eqref{am-riesz}, and put
$D_v=\prod_\tau|\det X_{c,\tau,v}|^{2a_\tau}$.
Then for all $F\in\bigoplus_cV$,
\begin{align}
 \|\Phi_c(F)\|_{\mathcal A}^2
 &\le D_v\left(\prod_{\tau,i}
                   s_{\tau,i}(R_{\tau,v})^{2\nu_{\tau,i}}\right)\|F\|_v^2
                          \label{am-class-sharp}\\
 &\le D_v\left(\prod_\tau
       \left(\sum_j\|y_{\tau,j,v}\|_{h_{c,\tau,v}}^2\right)^{d_\tau}\right)
                                                        \|F\|_v^2.
                          \label{am-class-hs}
\end{align}
The determinant contribution satisfies the exact identity
\begin{equation}\label{am-det-cancel}
 -\frac{\log N\mathfrak I_c}{d_K}
       +\frac1{2d_K}\sum_v n_v\log D_v
             =\sum_\tau a_\tau\widehat{\deg}(M_{c,\tau},h_{c,\tau}).
\end{equation}
\end{proposition}
\begin{proof}
Put $P_c=\rho_\lambda(X_{c,v}^{-1})F_c$ and
$d_v=\prod_\tau\det X_{c,\tau,v}^{a_\tau}$. The representation law gives
\[
             \rho_0(v(g))F_c=d_v\rho_0(R_v)P_c.
\]
Apply Theorem~\ref{am-singular} and
$\|P_c\|_V^2\le\sum_{c'}\|P_{c'}\|_V^2=\|F\|_v^2$.
Equation \eqref{am-riesz} proves the weaker bound.
Finally the exterior metric of the standard basis of $K^n$ is
$|\det X_{c,\tau,v}|^{-1}$. Definition \eqref{am-degree} therefore gives
\[
 \widehat{\deg}(M_{c,\tau},h_{c,\tau})
 =-\frac{\log N\mathfrak d(M_{c,\tau})}{d_K}
        +\frac1{d_K}\sum_v n_v\log|\det X_{c,\tau,v}|.
\]
Multiply by $a_\tau$ and sum. This also proves that the expression is
independent of all choices of coordinate bases and archimedean frames.
\end{proof}

\begin{corollary}[A bound from short integral functionals]\label{am-short}
Suppose
\[
 \widehat{\deg}(M_{c,\tau},h_{c,\tau})=-\tfrac12\log T
 \quad\text{for every $\tau$},\qquad T>0,
\]
and, for some $Q>0$,
\begin{equation}\label{am-short-assumption}
        \frac1{d_K}\sum_v n_v\|y_{\tau,j,v}\|_{h_{c,\tau,v}}^2
                                      \le Q\quad(\tau,j).
\end{equation}
Then the $c$th summand, with the metric \eqref{am-class-metric}, satisfies
\begin{equation}\label{am-final-slope}
 2\mu_{\max}(\Delta_\lambda(M_c))
 \le-\log w_{\min}+\sum_\tau d_\tau\log(nQ)
                                      -\sum_\tau a_\tau\log T.
\end{equation}
In particular, if $Q=n^2(2v_0)^2$ and
$c_{\nu}=C_\nu/\prod_i\nu_i!$, then
\begin{equation}\label{am-source-slope}
 2\mu_{\max}(\Delta_\lambda(M_c))
 \le\sum_\tau\left(
       d_\tau\log(4n^4v_0^2)+\log c_{\nu_\tau}-a_\tau\log T\right).
\end{equation}
One may further use
$c_\nu\le\prod_i\binom{\nu_i+n-i}{n-i}
\le(\nu_1+n)^{n(n-1)/2}$.
\end{corollary}
\begin{proof}
All $\|y_{\tau,j,v}\|^2$ are positive because the functionals form a
basis. The concavity bound $\log x\le x-1$, applied after division by
the mean, gives
\[
 \frac1{d_K}\sum_v n_v\log\left(\sum_j\|y_{\tau,j,v}\|^2\right)
 \le\log\left(\frac1{d_K}\sum_v n_v\sum_j\|y_{\tau,j,v}\|^2\right)
 \le\log(nQ).
\]
Use the integral image inclusion, Lemma~\ref{am-diagonal-slope},
\eqref{am-class-hs}, and the exact cancellation \eqref{am-det-cancel}.
This proves \eqref{am-final-slope}.
The multinomial bound preceding \eqref{am-weight-lower} gives
$w_{\min}^{-1}\le n^{\sum_\tau d_\tau}\prod_\tau c_{\nu_\tau}$.
For $Q=n^2(2v_0)^2$, combine the two logarithms of $n$ to obtain
\eqref{am-source-slope}. The bound on $c_\nu$ is
Theorem~\ref{am-gaussian}.
\end{proof}

\subsection{Base change, counterexamples, and the precise remaining input}

The data in the preceding corollary have their usual arithmetic meaning.
For a positive definite Hermitian space over a CM field $E$, integral
vectors in the Hermitian dual of a lattice give integral linear
functionals $z\mapsto h(y,z)$. A linearly independent $n$-tuple gives
an invertible matrix $g$, since nondegeneracy identifies a basis with
a basis of the dual. Upon an embedding $E\hookrightarrow K$, these
functionals extend to the $\mathcal O_K$-span of the lattice and remain
integral. Positivity and the Riesz identities persist at all embeddings.
No unimodularity of the selected $n$-tuple is required.

Normalized degree is unchanged by this base change. Indeed the
determinant ideal becomes $\mathfrak d(M)\mathcal O_K$, whose norm is
$N_E\mathfrak d(M)^{[K:E]}$; this follows prime by prime from
$\sum_{\mathfrak P\mid\mathfrak p}e_{\mathfrak P}f_{\mathfrak P}=[K:E]$.
Each archimedean term is repeated with total multiplicity $[K:E]$.
Both repetitions cancel the change from $[E:\mathbb Q]$ to
$[K:\mathbb Q]$ in \eqref{am-degree}. The same argument shows that
the mean of the squared norms in \eqref{am-short-assumption} is unchanged.
Thus a trace bound over the original CM field gives precisely the mean
bound used above, with no extra factor from the coefficient field.

Several possible simplifications are false, although none is used in
the actual coefficient-map statement.

\begin{itemize}
\item Negative determinant weights cannot be discarded. For $n=1$,
 $M=p\mathbb Z$, $a=-1$, and $\nu=0$, the transported lattice is
 $p^{-1}\mathbb Z$ and the coefficient map is the identity on constants.
 Its image is $p^{-1}\mathbb Z$, not $\mathbb Z$.
\item Integral functionals are essential for the determinant-ideal
 bound. For $M=\mathbb Z_p^2$, $a=0$, $\nu=(1,0)$, and
 $g=\operatorname{diag}(p^{-1},1)$, the polynomial $x_{11}$ acquires
 a coefficient $p^{-1}$. Formula \eqref{am-sharp-local}, which allows
 negative elementary divisors, still gives the correct answer.
\item The weighted coordinate metric is essential. For $\nu=(2,0)$
 in two rows, $C_\nu=2$ and the monomial $x_{11}x_{21}$ has normalized
 Fock norm $1/\sqrt2$. Its unweighted coefficient vector has norm one.
 Thus even $g=1$ would not give a contraction for the unweighted metric.
\item An integral independent $n$-tuple in a dual lattice need not be
 an integral basis. For $M=\mathbb Z_p^2$, the rows $(p,0),(0,1)$
 are independent and integral, but have index $p$ in the dual.
 The proof above uses only the inclusion $gM\subset\mathcal O^n$.
\item Stabilizers cannot be ignored when passing to class values. If
 a group of order two fixes the sole class representative and acts as
 $-1$ on a one-dimensional coefficient space in characteristic zero,
 every equivariant value is zero. It is not a full copy of the
 coefficient space. The actual full-product construction retains the
 trivial-stabilizer hypothesis.
\end{itemize}

The actual application supplies a decomposition into polynomial class
components, the local description \eqref{am-class-local}, and the
metric identity \eqref{am-class-metric}. Given these identifications,
this section proves the coefficient maps, their integral image,
their norm bounds, the determinant-frame cancellation, and the complete
deduction of \eqref{am-source-slope}. The existence of the decorated
base lattices, their specified determinant degree, and short independent
vectors in their Hermitian duals remain separate arithmetic inputs.
The estimates for the iterated period maps and their local vanishing
orders are also separate. In particular, the elementary-divisor formula
proved here must not be substituted for a valuation inequality for a
different map merely because both maps arise in the same construction.


